\subsection{由滤过定义的拓扑}

设 \(\mathbf{G}\) 为一个群，其滤过由 \((\mathbf{G}_n)_{n\in \mathbf{Z}}\) 这个 \(\mathbf{G}\) 的正规子群族给出。存在唯一一个与群结构相容的 \(\mathbf{G}\) 上的拓扑，使 \(G_{n}\) 构成 G 的单位元 e 的邻域基（《一般拓扑学》，第 III 章，§ 1，第 2 节，例）；称此为由滤过（\(G_{n}\)）在 G 上定义的拓扑。谈及滤过群的拓扑概念时，除非另有说明，均指由该滤过定义的拓扑。注意，\(G_{n}\) 是 G 的子群，因而既开又闭（《一般拓扑学》，第 III 章，§ 2，第 1 节，命题 4 的推论）。

由于每个 \(G_{n}\) 都是 G 的正规子群，G 上左一致结构和右一致结构的邻域相同；由此可知 G 有一个 Hausdorff 完备群 \(\hat{G}\)（《一般拓扑学》，第 III 章，§ 3，第 4 节，定理 1 及第 1 节，命题 2）。

对于 G 的每个子集 M，M 在 G 中的闭包等于

\[
\bigcap_ {n \in \mathbf {Z}} (\mathbf {M}. \mathbf {G} _ {n}) = \bigcap_ {n \in \mathbf {Z}} (\mathbf {G} _ {n}. \mathbf {M})
\]

	（《一般拓扑学》，第 III 章，§ 3，第 1 节，公式 (1)）；特别地，\(\bigcap_{n\in\mathbf{Z}}\mathrm{G}_{n}\) 是 \(\{e\}\) 的闭包；由此可见，\(\mathbf{G}\) 上的拓扑为 Hausdorff 拓扑的充要条件是滤过 \((\mathrm{G}_{n})\) 分离。\(\mathbf{G}\) 上的拓扑为离散拓扑的充要条件是存在 \(n\in\mathbf{Z}\) 使 \(\mathrm{G}_{n}=\{e\}\)（此时 \(\mathrm{G}_{m}=\{e\}\) 对 \(m\geqslant n\)）；这时称滤过 \((\mathrm{G}_{n})\) 为离散滤过。

由于与 G 相伴的 Hausdorff 群是 \(H = \frac{G}{\left(\bigcap_{n \in \mathbf{Z}} G_n\right)}\)，所以与之相伴的分次群 \(\text{gr}(G)\) 和 \(\text{gr}(H)\)（若给 H 赋予商滤过）可典范地识别。

现在设 \(G'\) 为另一个滤过群，\(u: G \to G'\) 为与滤过相容的同态；由 G 和 \(G'\) 上拓扑的定义立即可知 u 连续（*）。若 H 是 G 的子群（分别为正规子群），则 G 上拓扑在 H 上诱导的拓扑（分别为关于 H 的商拓扑）就是由 G 上滤过诱导的滤过在 H（分别 G/H）上定义的拓扑。G 和 \(G'\) 上拓扑的乘积拓扑，就是由 G 和 \(G'\) 上滤过的乘积定义的拓扑。

	令 \(v\) 为 \(G\) 上的阶函数（第 2 节）。关于 \(G_n\) 的假设意味着 \(v(xyx^{-1}) = v(y)\)，从而 \(v(xy^{-1}) = v(yx^{-1}) = v(x^{-1}y) = v(y^{-1}x)\)；对所有 \(x, y\) 属于 \(G\)，上述关系成立。取实数 \(\rho\)，使 \(0 < \rho < 1\)（例如取 \(\rho = 1 / e)\)），并令 \(d(x,y) = \rho^{y(xy - 1)}\)（对所有 \(x,y\) 属于 G）。则 \(d(x,x) = 0, d(x,y) = d(y,x)\)，而第 2 节的不等式 (5') 给出

\[
d (x, y) \leqslant \sup (d (x, z), d (y, z))\tag{15}
\]

对 G 中所有 x、y、z 都成立，从而得到三角不等式

\[
d (x, y) \leqslant d (x, z) + d (y, z).
\]

	因此 \(d\) 是 G 上在左、右平移下均不变的拟度量，而 \(\mathbf{G}_n\) 是由满足 \(x \in \mathbf{G}\) 且 \(d(e, x) \leqslant \rho^n\) 的元素构成的集合；由 \(d\) 定义的一致结构就是拓扑群 G 上的一致结构。若 G 是 Hausdorff 的，则 G 是零维可度量化拓扑空间（《一般拓扑学》，第 IX 章，§ 6，第 4 节）；并且，\(d\) 是 G 上的距离（只要滤过 \((\mathbf{G}_n)\) 也是穷竭的）。

回忆一下，对于拓扑环 A，左拓扑 A-模是一个 A-模 E，它带有与其加法群结构相容的拓扑，并且映射 \((a, x) \mapsto ax\) 从 A × E 到 E 连续（《一般拓扑学》，第 III 章，§ 6，第 6 节）。

命题 3. 设 A 为滤过环，\((\mathbf{A}_n)\) 为其滤过，B 为 A 的子环 \(\bigcup_{n\in \mathbf{Z}}\mathbf{A}_n\)，E 为滤过 B-模，\((\mathrm{E}_n)\) 为其滤过，F 为 E 的子 B-模 \(\bigcup_{n\in \mathbf{Z}}\mathrm{E}_n\)。则映射 \((a,x)\mapsto ax\) 从 \(\mathbf{B}\times \mathbf{F}\) 到 F 连续。

	设 \(a_0 \in \mathbf{B}\)、\(x_0 \in \mathbf{F}\)；根据假设，存在整数 \(r, s\) 使得 \(a_0 \in \mathbf{A}_r\) 且 \(x_0 \in \mathbf{E}_s\)。关系式

\[
a x - a _ {0} x _ {0} = (a - a _ {0}) x _ {0} + a _ {0} (x - x _ {0}) + (a - a _ {0}) (x - x _ {0})
\]

表明，若 \(a - a_0 \in \mathbf{A}_l\) 且 \(x - x_0 \in \mathrm{E}_f\)，则 \(ax - a_0x_0\) 属于

\[
\mathbf {E} _ {i + s} + \mathbf {E} _ {j + r} + \mathbf {E} _ {i + j}.
\]

	于是，给定整数 \(n\)，\(ax - a_0x_0 \in \mathrm{E}_n\) 只要 \(i \geqslant n - s\)、\(j \geqslant n - r\) 且 \(i + j \geqslant n\)，即只要 \(i\) 和 \(j\) 足够大。

推论。环 B 是拓扑环，B-模 F 是拓扑 B-模。

第一个断言由将命题 3 应用于 \(\mathbf{F} = \mathbf{B}_s\) 得到。

特别可见，滤过穷竭的滤过环 A 是拓扑环；在此情形下，任何滤过也穷竭的滤过 A-模都是拓扑 A-模。

	命题 4. 设 A 为由穷竭滤过 \((\mathbf{A}_n)\) 滤过的交换环，\(\mathfrak{p}\) 为 A 的一个理想。假设 \(\operatorname{gr}(\mathfrak{p}) = \bigoplus_{n \in \mathbf{Z}} (\mathfrak{p} \cap \mathbf{A}_n) / (\mathfrak{p} \cap \mathbf{A}_{n+1})\) 是环 \(\operatorname{gr}(\mathbf{A})\) 的素理想。则 \(\mathfrak{p}\) 在 A 中的闭包是素理想。

	我们知道 \(\operatorname{gr}(\mathbf{A} / \mathfrak{p})\) 同构于 \(\operatorname{gr}(\mathbf{A}) / \operatorname{gr}(\mathfrak{p})\)（第 4 节，命题 2），因而是整环；于是可知 \(\mathbf{A} / \bigcap_{n\in \mathbf{Z}}(\mathfrak{p} + \mathbf{A}_n)\) 是整环（第 3 节，命题 1 的推论）。因此，\(\bigcap_{n\in \mathbf{Z}}(\mathfrak{p} + \mathbf{A}_n)\) 是 \(\mathfrak{p}\) 的闭包，也是素理想。

设 A 为环，m 为 A 的双边理想；由 m-进滤过（第 1 节，例 3）在 A 上定义的拓扑称为 m-进拓扑；由于 m-进滤过是穷竭的，A 在此拓扑下是拓扑环（命题 3 的推论）。类似地，对每个 A-模 E，由 m-进滤过定义的拓扑称为 E 上的 m-进拓扑；E 在此拓扑下是拓扑 A-模。

	设 \(m'\) 是 A 的另一个双边理想；A 上的 \(m'\)-进拓扑比 \(m\)-进拓扑更细的充要条件是存在整数 n \textgreater{} 0 使得 \(m'^n \subset m\)；该条件是必要的，并且若它成立，则有 \(m'^{hn} \subset m^h\) 对所有 h \textgreater{} 0 成立，故该条件也是充分的。若 A 是交换 Noether 环，则这又等价于在 A 的素谱中有 \(V(m) \subset V(m')\)（第 II 章，§ 4，第 3 节，命题 11 的推论 2，以及 § 2，第 6 节，命题 15）。

